\chapter{欧氏空间中的微分形式}
\label{chap:differential-forms-euclidean}

上一章的三个积分定理使用了不同的符号：线积分、通量、散度和旋度。微分
形式把它们写成同一条公式
\[
 \int_Md\omega=\int_{\partial M}\omega.
\]
为了使这个公式有确定的含义，需要说明被积对象怎样随参数改变、边界怎样
定向。行列式的交错性给出前一种变换规则，外微分则把函数的导数、向量场
的旋度和散度放在相邻的次数中。本章先在线性空间中建立这些运算，再对
有限个正则参数块证明积分公式。

\section{交错形式与外代数}

设 $V$ 为有限维实向量空间。$k$-线性泛函
$\omega:V^k\to\R$ 若交换任意两个变量便变号，称为交错 $k$-形式；其空间
记为 $\Lambda^kV^*$。约定 $\Lambda^0V^*=\R$。

取基 $e_1,\ldots,e_n$ 及其对偶基 $e^1,\ldots,e^n$。对递增指标组
$I=(i_1,\ldots,i_k)$，定义
\[
 E_I(v_1,\ldots,v_k)=\det\bigl(e^{i_p}(v_q)\bigr)_{p,q=1}^k.
\]
它是交错形式。把每个向量按基展开，多重线性性表明任意交错形式由
其在基向量组上的值决定。重复指标的值为零，其余指标可排序，故
\[
 \omega=\sum_{i_1<\cdots<i_k}
 \omega(e_{i_1},\ldots,e_{i_k})E_I.
\]
在递增基向量组 $e_J$ 上，$E_I(e_J)$ 等于 $I=J$ 时的1或 $I\ne J$
时的0。因此这些形式构成一组基，维数为 $\binom nk$；$k>n$ 时空间为零。

对 $\alpha\in\Lambda^kV^*$、$\beta\in\Lambda^\ell V^*$，定义楔积
\[
 (\alpha\wedge\beta)(v_1,\ldots,v_{k+\ell})
 =\frac1{k!\ell!}\sum_{\sigma\in S_{k+\ell}}
 \operatorname{sgn}(\sigma)\,
 \alpha(v_{\sigma(1)},\ldots,v_{\sigma(k)})
 \beta(v_{\sigma(k+1)},\ldots,v_{\sigma(k+\ell)}).
\]
零次形式的楔积约定为标量乘法。置换和按前后两组内部的排列分组，每组
恰有 $k!\ell!$ 项；因两形式都交错，组内各项相同。于是对基形式，若
$I,J$ 有公共指标，$E_I\wedge E_J=0$；否则它等于连接指标组 $(I,J)$
排序后对应的基形式，乘以排序的符号。这也可在每个递增基向量组上直接
核验：只有把 $I$ 的指标送入第一组的项可能非零。

此乘法规则证明结合律：三组指标有重复时两种加括号方式都为零；无重复时
两者的符号都是把三组连接后排序的符号。由双线性性，结合律对任意形式
成立。特别地 $E_I=e^{i_1}\wedge\cdots\wedge e^{i_k}$。交换两组指标
需要 $k\ell$ 次对换，故
\[
 \alpha\wedge\beta=(-1)^{k\ell}\beta\wedge\alpha,
 \qquad \alpha\in\Lambda^kV^*,\ \beta\in\Lambda^\ell V^*.
\]
特别地，任一一次形式与自身的楔积为零；重复微分符号的项自动消失。

线性映射 $A:V\to W$ 的拉回定义为
\[
 (A^*\omega)(v_1,\ldots,v_k)=\omega(Av_1,\ldots,Av_k).
\]
若 $B:W\to Z$，则直接代入定义得到
$(B\circ A)^*=A^*\circ B^*$；拉回的次序与映射复合的次序相反。
把 $Av_j$ 代入楔积的置换和，还得到
$A^*(\alpha\wedge\beta)=A^*\alpha\wedge A^*\beta$。若 $V,W$ 同为
$n$ 维，分别选定坐标对偶基 $dx^i,dy^i$，则
$A^*(dy^1\wedge\cdots\wedge dy^n)=\det(A)dx^1\wedge\cdots\wedge dx^n$。
这里的行列式相对于所选基计算；它保留方向的符号。

\section{微分形式与外微分}

开集 $U\subset\R^n$ 上的微分 $k$-形式，是在每个点 $x$ 指定一个
交错形式 $\omega_x\in\Lambda^k(\R^n)^*$。记 $dx^i(v)=v_i$，
$dx^I=dx^{i_1}\wedge\cdots\wedge dx^{i_k}$，其中 $I$ 总按递增次序排列。
所谓 $C^r$，指下式的全部系数 $a_I$ 属于 $C^r(U)$：
\[
 \omega=\sum_{|I|=k}a_I(x)\,dx^{i_1}\wedge\cdots\wedge dx^{i_k}.
\]
光滑即 $C^\infty$ 形式的空间记为 $\Omega^k(U)$；零次形式就是函数。
对 $C^1$ 形式，外微分定义为
\[
 d\omega=\sum_{|I|=k}\sum_{j=1}^n
 \partial_ja_I\,dx^j\wedge dx^{i_1}\wedge\cdots\wedge dx^{i_k}.
\]

还需要把形式沿非线性映射拉回。若 $F:U\to W$ 是开集之间的 $C^1$
映射，$\omega$ 是 $W$ 上的连续 $k$-形式，定义
\[
 (F^*\omega)_x(v_1,\ldots,v_k)
 =\omega_{F(x)}(DF(x)v_1,\ldots,DF(x)v_k).
\]
零次时规定 $F^*f=f\circ F$。于是
\[
 F^*(a(y)\,dy^I)=(a\circ F)\,dF^{i_1}\wedge\cdots\wedge dF^{i_k}.
\]
链式法则给出 $(G\circ F)^*=F^*G^*$；逐点应用线性拉回的性质，得到
拉回与楔积相容。$F$ 为 $C^1$ 时拉回连续形式仍连续；若要对拉回再取
外微分，本章使用 $F\in C^2$、$\omega\in C^1$，使拉回系数确实为 $C^1$。

\begin{theorem}[外微分的代数性质]
\label{thm:exterior-derivative-properties}
若 $\alpha$ 是 $C^1$ 的 $k$-形式，$\beta$ 是 $C^1$ 形式，则
\begin{align*}
 d(\alpha\wedge\beta)&=d\alpha\wedge\beta+(-1)^k\alpha\wedge d\beta,
 &&\deg\alpha=k.
\end{align*}
若 $\alpha$ 为 $C^2$，则 $d(d\alpha)=0$。
若 $F:U\to W$ 为 $C^2$ 映射，$\omega$ 是 $W$ 上的 $C^1$ 形式，则
$d(F^*\omega)=F^*(d\omega)$。
\end{theorem}

\begin{proof}
只需先算单项。对 $\alpha=a\,dx^I$、$\beta=b\,dx^J$，
$d(ab)=b\,da+a\,db$。把 $db$ 越过 $k$ 个一次形式产生 $(-1)^k$，
便得第一式，再对有限和使用线性性。另一方面，
\[
 d^2(a\,dx^I)=\sum_{p,q}\partial_p\partial_q a\,
 dx^p\wedge dx^q\wedge dx^I=0.
\]
当 $p=q$ 时楔积为零，$p\ne q$ 时两项因 Schwarz 定理和交错性相消。

最后设 $\omega=a\,dy^I$。由 $d(dF^i)=0$ 及第一式，
\begin{align*}
 d(F^*\omega)
 &=d(a\circ F)\wedge dF^{i_1}\wedge\cdots\wedge dF^{i_k}\\
 &=\sum_j(\partial_j a\circ F)dF^j\wedge dF^{i_1}\wedge\cdots\wedge dF^{i_k}
 =F^*(d\omega).
\end{align*}
这一步只用到 $a\in C^1$ 和 $F\in C^2$。零次形式的等式就是链式法则。
\end{proof}

在 $\R^3$ 中，若 $\alpha=P\,dx+Q\,dy+R\,dz$，则
\[
 d\alpha
 =(Q_x-P_y)\,dx\wedge dy
 +(R_x-P_z)\,dx\wedge dz
 +(R_y-Q_z)\,dy\wedge dz.
\]
看似重复的九个偏导项中，有三个因 $dx^i\wedge dx^i=0$ 消失，其余六个
按反对称性两两合并。这个计算解释了旋度各分量中的符号从何而来。

对光滑形式，满足 $d\omega=0$ 的形式称为闭形式；能写成
$\omega=d\eta$、其中 $\eta$ 光滑的形式称为恰当形式。
$d^2=0$ 说明恰当必闭，反向是否成立取决于定义域。

\section{形式的积分与定向}

先取有界闭 Jordan 区域 $D\subset\R^k$。设 $X$ 在 $D$ 的一个开邻域
上为 $C^1$，在 $D$ 上单射且微分处处秩为 $k$。这样的参数片采用
$(\partial_1X,\ldots,\partial_kX)$ 给出的定向：两组有序切向基同向，
指它们之间的换基行列式为正。若 $\omega$ 在 $X(D)$ 的邻域连续，定义
\[
 \int_{X(D)}\omega:=\int_DX^*\omega.
\]
右端指 $X^*\omega=b(t)dt^1\wedge\cdots\wedge dt^k$ 的系数积分
$\int_D b(t)\,dt$，由连续性与 Jordan 条件可积。参数换元取闭域邻域上的
$C^1$ 微分同胚，映满相应参数域。
若重新参数化保持定向，拉回的复合律和变量代换中的 Jacobian 绝对值使积分
不变；反向定向则变号。有限参数片的内部不能重叠，公共边界须在各自参数
域中为 Jordan 零测集，随后按片相加。这里积分计的是定向区域，不把同一
片的多次覆盖自动当作一次；允许重覆盖时应明确另定义带重数的参数积分。

边界定向采用“先外法向，后边界正向基给出内部正向基”的约定。对平面区域，
这给出熟悉的逆时针外边界；对定向曲面，则等价于右手规则。
有向零维点的积分规定为带符号的函数值。因此正向区间 $[a,b]$ 的边界
是右端点减左端点，$\int_{\partial[a,b]}f=f(b)-f(a)$。

\section{Stokes 公式}

\begin{theorem}[欧氏参数化 Stokes 公式]
\label{thm:stokes-differential-forms-euclidean}
设 $1\le k\le n$，定向区域 $M\subset\R^n$ 是有限个参数块
$X_j(Q_j)$ 之并。每个 $Q_j$ 是非退化闭 $k$ 维矩形；$X_j$ 在其开邻域
为 $C^2$，在 $Q_j$ 上单射且处处秩为 $k$，并保持该片定向。
各块内部不交，块间只沿整张面或更低维的面相接。每张内部的
$(k-1)$ 维面恰为两块共有，两种参数由面邻域上的 $C^2$ 微分同胚联系，
诱导定向相反；其余 $(k-1)$ 维面的有向和记为 $\partial M$。
若 $\omega$ 是 $M$ 邻域中的 $C^1$ 的 $(k-1)$-形式，则
\[
 \int_Md\omega=\int_{\partial M}\omega.
\]
能细分为上述块的参数区域也适用。此处不主张任意光滑边界区域都已得到
这种有限分解；一般流形上的版本留待后文建立。
\end{theorem}

\begin{proof}
先对一个块 $M=X(Q)$ 证明，写 $Q=\prod_{j=1}^k[\ell_j,u_j]$。
参数 $X$ 为 $C^2$，故 $X^*\omega$ 为
$C^1$，且 $X^*(d\omega)=d(X^*\omega)$。写
\[
 X^*\omega=\sum_{j=1}^k(-1)^{j-1}a_j(t)\,
 dt^1\wedge\cdots\wedge\widehat{dt^j}\wedge\cdots\wedge dt^k.
\]
于是
\[
 d(X^*\omega)=\left(\sum_{j=1}^k\partial_ja_j\right)
 dt^1\wedge\cdots\wedge dt^k.
\]
对第 $j$ 项沿 $t^j$ 使用 Newton--Leibniz 公式，其余坐标使用 Fubini
定理，得到矩形上、下两张 $j$-边界面的积分；系数 $(-1)^{j-1}$ 正是诱导
边界定向的符号。

具体地，上面 $t^j=u_j$ 的定向相对于删去 $e_j$ 的坐标基具有符号
$(-1)^{j-1}$，下面 $t^j=\ell_j$ 的符号相反，因而两面的贡献恰为
$a_j(\ldots,u_j,\ldots)-a_j(\ldots,\ell_j,\ldots)$ 的其余坐标积分。
当 $k=1$ 时这就是两端点之差。

对所有块相加。公共面的两种参数表示由换元公式给出相同的绝对积分，
定向相反使它们抵消。更低维的面在面参数中为 Jordan 零测集；这一点也可
用低维 Lipschitz 像的矩形覆盖估计验证，所以不会另生积分项。剩下的
正是按所定方向计算的 $\partial M$ 积分。
\end{proof}

\section{星形域上的 Poincaré 引理}

开集 $U$ 关于 $a$ 星形，是指每条线段 $[a,x]$ 都包含于 $U$。

\begin{theorem}[Poincaré 引理]
\label{thm:poincare-lemma-star-shaped}
若 $U\subset\R^n$ 关于 $0$ 星形，$k\ge1$，且
$\omega\in\Omega^k(U)$ 闭，则存在 $\eta\in\Omega^{k-1}(U)$ 使
$\omega=d\eta$。
\end{theorem}

\begin{proof}
令径向向量场 $R_x=x$，定义同伦算子
\[
 (H\omega)_x(v_1,\ldots,v_{k-1})
 =\int_0^1t^{k-1}\omega_{tx}(x,v_1,\ldots,v_{k-1})\,dt.
\]
先解释求导为什么能进入积分。固定 $x_0\in U$，线段 $[0,x_0]$ 紧且
包含于开集 $U$，故它的某个闭管状邻域仍在 $U$ 内。$x$ 足够接近 $x_0$
时全部线段 $[0,x]$ 都落在这个共同紧集内。被积系数的各阶导数在那里
连续，差商由一元中值公式一致趋于相应偏导，所以可逐次移入积分，
$H\omega$ 因而光滑。

下面完整计算一个单项 $\omega=a(x)\theta$，其中 $\theta=dx^I$ 为
常系数 $k$-形式。用插入径向向量的记号
\[
 (\iota_R\theta)_x(v_1,\ldots,v_{k-1})=\theta(x,v_1,\ldots,v_{k-1}).
\]
按第一列展开行列式，得到
\[
 \iota_R\theta=\sum_{r=1}^k(-1)^{r-1}x_{i_r}
 dx^{i_1}\wedge\cdots\wedge\widehat{dx^{i_r}}\wedge\cdots\wedge dx^{i_k},
 \qquad d\iota_R\theta=k\theta.
\]
帽号表示删去该因子；求导后把 $dx^{i_r}$ 移回原位，产生的符号恰好
抵消前面的 $(-1)^{r-1}$，所以有 $k$ 个 $\theta$。
对一次形式 $\lambda$，由交错展开还有
$\iota_R(\lambda\wedge\theta)=\lambda(R)\theta-\lambda\wedge\iota_R\theta$。
记 $b_t(x)=a(tx)$、$\lambda_t=\sum_j\partial_j a(tx)dx^j$，于是
\begin{align*}
 dH\omega&=\int_0^1\left(t^k\lambda_t\wedge\iota_R\theta
                       +k t^{k-1}b_t\theta\right)dt,\\
 Hd\omega&=\int_0^1t^k\left(\sum_jx_j\partial_ja(tx)\theta
                       -\lambda_t\wedge\iota_R\theta\right)dt.
\end{align*}
相加后含 $\lambda_t$ 的两项抵消，余下为
$\int_0^1\frac{d}{dt}(t^ka(tx))\,dt\,\theta= a(x)\theta$。
对所有单项相加得到同伦恒等式
\[
 dH\omega+Hd\omega=\omega-s_0^*\omega,
\]
其中 $s_0(x)=0$；右端两项分别来自 $t=1$ 与 $t=0$。当 $k\ge1$ 时，
常值映射的导数为零，所以 $s_0^*\omega=0$。又因 $d\omega=0$，
$d(H\omega)=\omega$。令 $\eta=H\omega$ 即得。
零次时约定 $Hf=0$，直接由一元基本定理有
$Hdf=\int_0^1\sum_jx_j\partial_jf(tx)dt=f(x)-f(0)$，所以同伦恒等式
在零次也成立，但闭零次形式不在引理的恰当性结论之内。
\end{proof}

对任意开集 $U$ 和任意 $x\in U$，都可取一个以 $x$ 为中心、闭包仍包含于
$U$ 的小开球。开球关于其中心星形，因而上一定理逐点给出：闭形式总是局部
恰当。要把这些局部原函数拼成全局原函数，还需要控制定义域的拓扑；穿孔平面
上的角形式 $(-y\,dx+x\,dy)/(x^2+y^2)$ 就是不能全局拼合的基本例子。

\section{经典公式的形式表述}

在 $\R^3$ 中，标量函数 $f$ 对应 $0$-形式，向量场
$F=(P,Q,R)$ 可对应
\[
 \alpha=Pdx+Qdy+Rdz,
 \qquad
 \beta=P\,dy\wedge dz+Q\,dz\wedge dx+R\,dx\wedge dy.
\]
于是 $df$ 对应梯度，$d\alpha$ 对应旋度，$d\beta$ 对应散度。对曲线、曲面
和区域应用同一 Stokes 公式，分别得到 Newton--Leibniz、经典 Stokes 与
Gauss 定理的公式；二维的 $Pdx+Qdy$ 则给出 Green 公式。
作为定理的推导，这里限于已经满足本章有限正则分块条件的区域；记号上的
对应不能代替对一般积分域的证明。

这种对应依赖欧氏内积和定向，而外微分本身不依赖度量。后续流形章节会把
积分与 Stokes 定理保留在抽象流形上；梯度、散度和 Hodge 星算子则需要额外
的度量结构。

\section*{习题}
\addcontentsline{toc}{section}{习题}
\subsection*{A组　外代数与外微分}
\begin{enumerate}[label=\textbf{\thechapter.\arabic*.}]
 \exerciseitem{alternating-form-basis}
 证明 $\dim\Lambda^k(\R^n)^*=\binom nk$，并写出 $n=4,k=2$ 的标准基。
 \exerciseitem{wedge-sign-computation}
 计算 $(dx+dy)\wedge(dx-dy)$ 及
 $(dx+dy)\wedge(dy+dz)\wedge(dz+dx)$，说明含重复一次形式因子的项为何为零。
 \exerciseitem{top-form-determinant-pullback}
 证明线性映射对顶次体积形式的拉回乘以行列式。
 \exerciseitem{exterior-derivative-leibniz}
 在坐标中证明外微分的分次 Leibniz 公式。
 \exerciseitem{exterior-derivative-square-zero}
 完整证明 $d^2=0$，指出 Schwarz 定理和交错性各自的作用。
 \exerciseitem{pullback-commutes-d}
 对一次形式先证明 $d(F^*\omega)=F^*(d\omega)$，再推广到一般形式。
\end{enumerate}
\subsection*{B组　Stokes 与 Poincaré}
\begin{enumerate}[label=\textbf{\thechapter.\arabic*.},resume]
 \exerciseitem{stokes-rectangle-coordinate-proof}
 在三维长方体上对一般二次形式逐项验证 Stokes 公式。
 \exerciseitem{boundary-orientation-square}
 从“外法向优先”约定推导平面正方形边界的逆时针方向。
 \exerciseitem{poincare-one-form-potential}
 把同伦算子用于闭一次形式，写出显式势函数并直接求导验证。
 \exerciseitem{poincare-homotopy-identity}
 对一般 $k$-形式逐项验证 $dH+Hd=I-s_0^*$。
 \exerciseitem{punctured-plane-closed-not-exact-form}
 证明角形式闭但不恰当，并用单位圆积分说明原因。
 \exerciseitem{local-poincare-lemma}
 由星形域版本证明闭形式在每一点附近恰当。
\end{enumerate}
\subsection*{C组　经典对应}
\begin{enumerate}[label=\textbf{\thechapter.\arabic*.},resume]
 \exerciseitem{green-as-differential-form}
 把 Green 定理写成一次形式的 Stokes 公式并核对符号。
 \exerciseitem{gauss-as-differential-form}
 计算与向量场对应的二次形式的外微分，恢复散度定理。
 \exerciseitem{curl-as-exterior-derivative}
 计算 $d(Pdx+Qdy+Rdz)$ 并恢复旋度的三个分量。
 \optionalexerciseitem{forms-metric-independence-audit}
 区分本章哪些构造只需光滑结构与定向，哪些还使用欧氏内积。
\end{enumerate}
